\scanpage{57}
\mspage{9}
Notons que l'isom
\[
\psi_r:\Pi_{03}\xrightarrow{\sim}\Pi_\omega
\]
\[
\inserted{\begin{gathered}
r\in\operatorname{Rep}(J,\uline{\bar\omega})\quad\text{\deleted{\unsure{pour}}}\\
\simeq\operatorname{Isom}((\Sigma_0^*,P_+),(\Sigma^*,P_+))
\end{gathered}}
\]
est compatible avec le transport des structures,
\deleted{ce qui \unsure{fait}} \inserted{donc pour} $\forall$ élément
\[
\begin{array}{rll}
g\in{}&\widetilde G_J&{}\subset\Gamma=G_J\\
&\inserted{\Vert}&\\
&\inserted{\Gamma_{P_\varepsilon}}&
\inserted{\begin{aligned}[t]
&=\left\{(g,\tau^{\unsure{\operatorname{sg}g}})\in G_J\times \deleted{\unsure{\{1,\tau\}}}\inserted{\mathbb Z/2\mathbb Z}
\ \middle|\ \begin{array}{l}
\alpha=0\text{ si }\operatorname{sg}g=1\\
\alpha=1\text{ si }\operatorname{sg}g=-1
\end{array}\right\}\\
&=\{g\in\Gamma\mid g_\varepsilon^{(\varepsilon)}=\operatorname{id}_\varepsilon\}
\end{aligned}}
\end{array}
\]
on aura commutativité de
\diagram{
\begin{tabular}{@{}c@{\qquad}c@{}}
$(37)$ &
\begin{tikzcd}[ampersand replacement=\&,column sep=large,row sep=large]
\unsure{\downarrow|\vee}\&\Pi_{03}
\arrow[dl,"{\psi_r=\psi_{\omega,i}^\varepsilon}",description]
\arrow[dr,"{\psi_{g(r)}=\psi_{g(\omega),g(i)}^{g(\varepsilon)}}"]\&\\
\Pi_{\unsure{\varepsilon}}\arrow[rr,"{\Pi_1(g)}","{\begin{gathered}\text{(par transport}\\\text{de structure}\end{gathered}}"']\&\&\Pi_\varepsilon
\end{tikzcd}
\end{tabular}
}
ce qui suffit en principe à passer d'un $\psi_r$ à
$\psi_{g(r)}=\deleted{\unsure{\varphi|\psi}}\cdot\Pi_1(g)\circ\psi_r$, vu que \deleted{\unsure{g}}
$\widetilde G_J$ est simplement transitif sur l'ens. des repères $r$
correspondant à un $\varepsilon$ fixé.
Mais on aimerait exprimer le passage par une formule de transformation
\uline{dans} $\Pi_{0,3}$.

Notons que la donnée du repère \inserted{$r$} définit un isom canonique
(encore noté ${\unsure{u}}_r$)
\margin{56}
